\begin{question}
\noindent An e-sports arena has a rotating stage logo shaped like a regular hexagon, used as the map-select icon in a Valorant tournament broadcast. The graphic designer draws the hexagonal logo on a coordinate grid centred at the origin $O$, as shown below.

\begin{center}
\begin{tikzpicture}[scale=1.1]
    \coordinate (O) at (0,0);
    \def\R{3}

    % Regular hexagon vertices, one vertex pointing right along positive x-axis
    \coordinate (V1) at (0:\R);
    \coordinate (V2) at (60:\R);
    \coordinate (V3) at (120:\R);
    \coordinate (V4) at (180:\R);
    \coordinate (V5) at (240:\R);
    \coordinate (V6) at (300:\R);

    \draw[thick, fill=cyan!15] (V1) -- (V2) -- (V3) -- (V4) -- (V5) -- (V6) -- cycle;

    \fill[black] (O) circle (1.5pt);
    \node[below left] at (O) {$O$};

    \node[right] at (V1) {$A$};
    \node[above right] at (V2) {$B$};
    \node[above left] at (V3) {$C$};
    \node[left] at (V4) {$D$};
    \node[below left] at (V5) {$E$};
    \node[below right] at (V6) {$F$};

    % dashed lines from centre to two adjacent vertices to show rotation angle
    \draw[dashed] (O) -- (V1);
    \draw[dashed] (O) -- (V2);

    \pic [draw, angle radius=0.7cm, "$x^\circ$"] {angle = V1--O--V2};

\end{tikzpicture}
\end{center}

\noindent The logo $ABCDEF$ is a regular hexagon with centre $O$. During the broadcast animation, the logo is rotated repeatedly about $O$ so that it always lands back exactly on top of itself (vertex meeting vertex).

\begin{enumerate}
    \item[(a)] Work out the smallest angle $x$, marked on the diagram, through which the logo must be rotated about $O$ so that it maps onto itself.

    \vspace{2.5cm}
    \answerequnit{x}{$^\circ$}{1}

    \item[(b)] State the order of rotational symmetry of the hexagonal logo.

    \vspace{2cm}
    \answerplain{1}

    \item[(c)] The animation designer instead wants the logo to complete a sequence of rotations that brings vertex $A$ to the position currently occupied by vertex $D$, using the \textbf{smallest possible number of equal individual rotation steps} (each step being one of the symmetry rotations found in part (a)), and each step must be in the \textbf{same direction}. Work out the number of individual rotation steps needed, and state the direction (clockwise or anticlockwise) that requires the fewest steps.

    \vspace{3.5cm}
    \answerplain{1}
\end{enumerate}
\end{question}
